\documentclass[preview]{standalone}
\usepackage[english]{babel}
\usepackage{amsmath}
\usepackage{amssymb}
\begin{document}
\begin{center}
\special{dvisvgm:raw <g id='unique000'>}Dado un símplex (ej. un triángulo con vértices $P_0, P_1, P_2$),\\ cualquier punto $P$ se expresa de forma única como:\\$P = \lambda_0 P_0 + \lambda_1 P_1 + \lambda_2 P_2$\\donde $\lambda_0 + \lambda_1 + \lambda_2 = 1$.\\$(\lambda_0, \lambda_1, \lambda_2)$ son las \textbf{coordenadas baricéntricas} de $P$.\special{dvisvgm:raw </g>}
\end{center}
\end{document}